\chapter{Predicting Representation Need}
\label{ch:pred}

\chapterhandoff{Chapter~\ref{ch:repr} established a retrospective ceiling and
showed how much of it repetition already buys. The ceiling is computed with
hindsight. This chapter asks whether the same decision can be made
\emph{beforehand}, from signals a deployment would actually have, and is
designed so that the answer --- whichever way it falls --- cannot be attributed
to the protocol.}

\section{Which half of the problem to attack}
\label{sec:which-half}

Section~\ref{sec:router-formulation} split the oracle into triage (does any mode
solve this task?) and which-mode (given that one does, which?). They differ in
one respect that determines what is even testable.

The which-mode label is minted only by a solve event. Table~\ref{tab:supply}
shows the consequence: across the six VisualWebArena cells the number of
\emph{minted} which-mode labels ranges from 15 to 97, pooled total 260, and after
requiring at least two classes each retaining enough members to survive a
five-fold split, \textbf{4 of 6 cells fail that admissibility rule entirely}.

\begin{scopebox}
\textbf{Say exactly what that rule is, because the stronger phrasing would be
false.} ``No trainable classifier'' is an \emph{admissibility} verdict under a
stated minimum of 10 training rows per class ($n\times\tfrac45 \geq 10$), not a
mathematical impossibility. A logistic regression will happily fit two classes
with three examples each and return a number; the rule exists to stop us
believing it. The defensible statement is therefore \emph{4 of 6 cells fail the
$\mathrm{min\_class\_n}=10$ rule}, and what it supports is a claim about how much
supervision the benchmark produced --- not about what an estimator can be made to
output. Two further caveats on the same table: the rule is evaluated against the
\emph{expected} fold size $n\times\tfrac45$ rather than realised fold membership,
so a cell can pass it and still lack a complete five-fold classifier; and the
counts are \emph{minted} labels, of which about 54\% pooled fall in classes that
survive the filter.
\end{scopebox}

\begin{table}[t]
\centering\small
\caption[Which-mode label supply]{\textbf{The which-mode half fails its
admissibility rule before modelling begins.} A label is \emph{minted} only where
some mode succeeded. The rule requires two classes each retaining $\geq10$
expected training rows under a five-fold split; four cells do not have them. This
is an admissibility verdict, not an impossibility proof (see the box above), and
Chapter~\ref{ch:disc} argues that the underlying constraint is supply.}
\label{tab:supply}
\begin{tabular}{@{}lrrrrl@{}}
\toprule
Cell & scored $n$ & minted labels & solvable \% & classes surviving & admissible? \\
\midrule
\cell{cls$\cdot$B0} & 224 & 97 & 43.30 & 3 (\mDOM{}, \mPP{}, \mSoM{}) & yes \\
\cell{red$\cdot$B0} & 203 & 53 & 26.11 & 1 (\mDOM{})                  & \textbf{no} \\
\cell{cls$\cdot$B1} & 224 & 55 & 24.55 & 2 (\mDOM{}, \mSoM{})         & yes \\
\cell{red$\cdot$B1} & 203 & 24 & 11.82 & 0                             & \textbf{no} \\
\cell{cls$\cdot$B2} & 224 & 16 & \phantom{0}7.14 & 0                   & \textbf{no} \\
\cell{red$\cdot$B2} & 203 & 15 & \phantom{0}7.39 & 0                   & \textbf{no} \\
\bottomrule
\end{tabular}
\end{table}

The triage label is better posed. It is binary, defined on \emph{every}
task, and available in all 8 cells including the WebArena ones, and it is the
half that Chapter~\ref{ch:repr} showed to be worth 9.5--30.6\% cost at zero
success loss under an oracle. \textbf{We therefore run the predictive study on
triage} --- deliberately the most favourable framing available: two classes
instead of six, a policy whose oracle version cannot lose success, and a
comparator that is trivially implementable. A negative result here is stronger
than a negative result on the harder half.

\begin{scopebox}
\textbf{But note what triage does \emph{not} escape.} Its positive class is the
solvable set --- precisely the tasks that mint which-mode labels. The counts in
the ``minted labels'' column of Table~\ref{tab:supply} are simultaneously the
which-mode label supply \emph{and} the triage minority-class size: 97, 53, 55,
24, 16, 15. Both halves are rationed by the same scarce resource; triage merely
divides it two ways instead of six. So triage is the easier problem, not a
problem with more evidence behind it --- and on \cell{cls$\cdot$B2} the
classifier is being asked to learn from 16 positives in 224 rows, roughly three
per fold.
\end{scopebox}

\section{The serving-time feature contract}
\label{sec:features}

A router feature is admissible only if a deployment would possess it at decision
time, before the expensive observation is constructed. Twenty raw features meet
that bar on VisualWebArena:

\begin{itemize}[leftmargin=1.5em]
\item \textbf{Five numeric}: \texttt{dom\_complexity}, \texttt{text\_length},
      \texttt{tokens\_input\_text}, \texttt{intent\_token\_count},
      \texttt{reasoning\_difficulty}.
\item \textbf{Fifteen binary}: \texttt{has\_reference\_image} plus fourteen
      lexical intent predicates matched against the task instruction, namely
      \texttt{intent\_}$\langle$\emph{search, compare, visual\_attribute,
      form\_fill, sort, filter, nav, temporal, color, aggregate, compose,
      question, action\_word, account\_action}$\rangle$. Their definitions are
      in Appendix~\ref{app:router}.
\end{itemize}

All are computable from the task instruction and the cheap structured
observation. None requires running a model, constructing an image, or observing
an outcome. Appendix~\ref{app:router} gives the complete dictionary with
per-feature definitions, together with the exact order of operations inside a
cross-validation fold, so that the leakage controls described below can be
checked rather than taken on trust.

\subsection{Two features do not exist on WebArena, and this matters twice}
\label{sec:matched-features}

WebArena carries no reference images and no difficulty annotation
(Figure~\ref{fig:f4}(B)). Zero-filling them there and comparing against the
20-feature VisualWebArena numbers would place an unmatched column beside matched
ones. We therefore drop \texttt{has\_reference\_image} and
\texttt{reasoning\_difficulty} on \emph{every} cell and refit all eight, giving
a matched 18-feature set. The 20-feature VisualWebArena results are reported
separately and are not a subset of the matched table --- every number was
recomputed.

\begin{table}[t]
\centering\small
\caption[Effect of the matched feature set]{\textbf{Dropping two features that
no deployment has costs most of the discriminative power.} AUROC for the triage
label, out-of-fold, per cell. Five of six VisualWebArena cells fall, one by
0.141.}
\label{tab:feature-ablation}
\begin{tabular}{@{}lrrr@{}}
\toprule
Cell & 20 features (VWA-only) & 18 features (matched) & $\Delta$ \\
\midrule
\cell{cls$\cdot$B0} & 0.726 & 0.683 & $-0.043$ \\
\cell{red$\cdot$B0} & 0.780 & 0.700 & $-0.080$ \\
\cell{cls$\cdot$B1} & 0.732 & 0.705 & $-0.027$ \\
\cell{red$\cdot$B1} & \textbf{0.864} & 0.723 & $\mathbf{-0.141}$ \\
\cell{cls$\cdot$B2} & 0.642 & 0.646 & $+0.004$ \\
\cell{red$\cdot$B2} & 0.615 & 0.526 & $-0.089$ \\
\bottomrule
\end{tabular}
\end{table}

\begin{scopebox}
\textbf{First, the order of events, because it matters for whether this is
leakage.} The two features were dropped because WebArena \emph{does not have
them} --- a property of the corpus, established by the EDA in
Chapter~\ref{ch:system} before any model was fitted, and visible in
Figure~\ref{fig:f4}(B). They were not dropped because they performed too well.
Had we inspected out-of-fold importances and removed whatever looked suspicious,
that would have been whole-dataset feature selection dressed as nesting, and the
nesting claim in Section~\ref{sec:protocol} would be void. The admissibility
rule is availability at serving time, applied to the feature set as a whole.

\textbf{What we then found is one of this chapter's main results.} In the
20-feature configuration the strongest single feature in 5 of 6
cells is \texttt{reasoning\_difficulty} --- \emph{VisualWebArena's own
human-authored annotation of how hard the task is}. A router reading it is
reading the benchmark's statement of the answer's difficulty, which no
deployment possesses. Removing it, and the equally unavailable reference-image
flag, is what causes the fall in Table~\ref{tab:feature-ablation}. This is
independent evidence for the mechanism argued in Chapter~\ref{ch:disc}, and it
is also the cheapest available rebuttal to the objection that our negative
result is merely undersampling: the issue is not that there are too few rows,
but that the one column reliably carrying signal does not exist outside the
benchmark.
\end{scopebox}

\begin{figure}[t]
  \centering
  \includegraphics[width=\textwidth]{fig_f11_feature_set_ablation.pdf}
  \caption[Two feature sets side by side]{%
  \textbf{What the router loses when it is restricted to features a deployment
  has.} Per-cell AUROC under the 20-feature VisualWebArena set and the matched
  18-feature set. The two dropped features are the benchmark's own difficulty
  annotation and its reference-image flag. Read together with
  Figure~\ref{fig:f4}(B), which shows why they had to be dropped: WebArena has
  neither.}
  \label{fig:f11}
\end{figure}

\subsection{A matched set is not a matched support}

Matching the feature \emph{set} does not match the feature \emph{support}. Five
of the fourteen intent regexes have zero variance on the WebArena cells: they
were written against VisualWebArena phrasing, and WebArena words its intents
differently. So the WebArena cells are effectively fitted on 13 live features
against the VisualWebArena cells' 15--17 (Table~\ref{tab:live-features}). We
state the comparison as matched on set, not on support.

\begin{table}[t]
\centering\small
\caption[Live features per cell]{\textbf{A column with no variance carries no
information.} The header figure of 18 is a ceiling, not a count of usable
features.}
\label{tab:live-features}
\begin{tabular}{@{}llp{7.4cm}@{}}
\toprule
Cell & live & dead (zero-variance) \\
\midrule
\cell{cls$\cdot$B0}, \cell{cls$\cdot$B1}, \cell{cls$\cdot$B2} & 15/18 &
  \texttt{intent\_account\_action}, \texttt{intent\_action\_word},
  \texttt{intent\_filter} \\
\cell{red$\cdot$B0}, \cell{red$\cdot$B1}, \cell{red$\cdot$B2} & 17/18 &
  \texttt{intent\_action\_word} \\
\cell{wa\_red$\cdot$B0}, \cell{wa\_red$\cdot$B1} & 13/18 &
  \texttt{intent\_action\_word}, \texttt{intent\_color},
  \texttt{intent\_compare}, \texttt{intent\_filter},
  \texttt{intent\_form\_fill} \\
\bottomrule
\end{tabular}
\end{table}

\section{Protocol: what a negative result has to survive}
\label{sec:protocol}

The router is an $L_2$-regularised logistic regression on the 18 raw features,
evaluated by task-held-out five-fold cross-validation per cell, seed 42. The
model is deliberately simple; the interesting design decisions are all in what
is held out.

\paragraph{Full nesting.} Within each outer fold: feature standardisation is fit
on training rows only; the best-SR mode and the cheapest mode are re-selected
from \emph{training-row} success and cost; the operating threshold $\tau$ is
chosen against inner-CV out-of-fold scores over training rows only; and the
outer-test rows are scored by a model fitted on training rows alone. Nothing
that touches an outer test fold has seen it.

This matters because we previously got it wrong. An earlier revision of this
analysis reported a ``nested'' operating point while reusing globally computed
out-of-fold scores --- whose folds include the outer test rows --- together with
a whole-cell choice of best and cheapest mode. The numbers below supersede it.
We record the correction because it changes the interpretation of a headline
number, not to narrate debugging.

\paragraph{One thing nesting exposed.} The best mode is \emph{not stable across
folds}. On \cell{red$\cdot$B0} the five outer folds select \mDOM{}, \mDOM{},
\mSoM{}, \mSoM{}, \mDOM{}. A pipeline that picks one best mode from all realised
outcomes is not merely optimistic about its threshold --- it reports a mode
choice its own resampling does not reproduce.

\paragraph{Two controls, neither optional.} \emph{(i)} A bundle-permutation
null at $B=10{,}000$: the permutation unit is the whole task bundle
$(y,\,\mathrm{success},\,\mathrm{cost})$ against $X$. Permuting $y$ alone would
disconnect the label from the outcomes that define it, and its error is not
one-directional. \emph{(ii)} The always-cheapest fixed policy as the comparator,
because it is the policy a deployment gets for free.

\paragraph{Baseline sanity check.} We report each cell's best \emph{single}
covariate alongside the fitted model. If the model does not clear it, ``learnable''
means ``learnable by one covariate'', which is not a router.

\section{Is the triage label predictable at all?}
\label{sec:auroc}

\begin{table}[t]
\centering\small
\caption[Triage predictability on the matched feature set]{\textbf{Out-of-fold
AUROC for the triage label, all eight cells, 18 matched features.} A
prior-only predictor scores 0.500 by construction. The model clears its own best
single covariate in 5 of 8 cells, by $+0.014$ to $+0.075$.}
\label{tab:auroc}
\begin{tabular}{@{}lrrrrl@{}}
\toprule
Cell & $n$ & solvable \% & AUROC (LR) & best single feature & that feature \\
\midrule
\cell{cls$\cdot$B0}     & 224 & 43.3 & \textbf{0.683} & 0.607 & \texttt{dom\_complexity} \\
\cell{red$\cdot$B0}     & 203 & 26.1 & \textbf{0.700} & 0.685 & \texttt{intent\_token\_count} \\
\cell{cls$\cdot$B1}     & 224 & 24.6 & \textbf{0.705} & 0.631 & \texttt{intent\_token\_count} \\
\cell{red$\cdot$B1}     & 203 & 11.8 & \textbf{0.723} & 0.664 & \texttt{intent\_token\_count} \\
\cell{cls$\cdot$B2}     & 224 & \phantom{0}7.1 & 0.646 & 0.655 & \texttt{intent\_compare} \\
\cell{red$\cdot$B2}     & 203 & \phantom{0}7.4 & 0.526 & \textbf{0.790} & \texttt{intent\_token\_count} \\
\cell{wa\_red$\cdot$B0} & 104 & 51.9 & 0.554 & 0.607 & \texttt{dom\_complexity} \\
\cell{wa\_red$\cdot$B1} & 104 & 30.8 & \textbf{0.758} & 0.710 & \texttt{dom\_complexity} \\
\bottomrule
\end{tabular}
\end{table}

The answer is a qualified yes, and the qualification is the interesting part.
AUROC spans 0.526--0.758. In five cells the fitted model beats its own best
single covariate; in three it does not, and in \cell{red$\cdot$B2} it is beaten
by 0.264 --- the single feature \texttt{intent\_token\_count} scores 0.790 while
the model scores 0.526, barely above chance.

\begin{scopebox}
\textbf{A retraction of our own earlier reading.} An earlier version of this
analysis summarised these results as ``AUROC 0.65--0.72 in five of six cells'',
which implies adequate discriminative power throughout. The sixth cell,
\cell{red$\cdot$B2}, scores 0.483--0.526 depending on feature set --- at or below
chance --- and it is precisely the cell whose saving survives significance
testing in Section~\ref{sec:permutation}. Global discrimination and tail
usefulness come apart on this data. The corrected statement is
Section~\ref{sec:auroc-not-answer}.
\end{scopebox}

\section{Is the saving real, or manufactured by the sweep?}
\label{sec:permutation}

The operating threshold is chosen post hoc by a sweep, so some apparent saving
is guaranteed even from a signal-free pipeline. The bundle-permutation null
measures how much.

\begin{table}[t]
\centering\small
\caption[Permutation control]{\textbf{How much of each cell's saving a
signal-free pipeline reproduces.} $B=10{,}000$ bundle permutations per cell;
$p$ is the plus-one Monte Carlo estimator $(k+1)/(B+1)$, smallest reportable
$1.00\times10^{-4}$. Holm's tightest threshold over $m=8$ cells is
$0.05/8 = 6.25\times10^{-3}$.}
\label{tab:permutation}
\begin{tabular}{@{}lrrr@{}}
\toprule
Cell & observed saving & median under shuffled labels & $p$ \\
\midrule
\cell{cls$\cdot$B0}     & \phantom{0}7.0\% & \phantom{0}3.3\% & 0.242 \\
\cell{red$\cdot$B0}     & \phantom{0}0.0\% & \phantom{0}0.0\% & 1.000 \\
\cell{cls$\cdot$B1}     & \phantom{0}2.6\% & 10.1\% & 0.741 \\
\cell{red$\cdot$B1}     & 11.8\% & \phantom{0}1.0\% & 0.033 \\
\cell{cls$\cdot$B2}     & 20.4\% & 20.8\% & 0.682 \\
\cell{red$\cdot$B2}     & 26.7\% & \phantom{0}5.5\% & \textbf{0.0004} \\
\cell{wa\_red$\cdot$B0} & \phantom{0}4.4\% & \phantom{0}0.6\% & 0.186 \\
\cell{wa\_red$\cdot$B1} & \phantom{0}3.8\% & \phantom{0}2.9\% & 0.443 \\
\bottomrule
\end{tabular}
\end{table}

Under Holm's step-down procedure \citep{holm1979} at $\alpha=0.05$ over the
$m=8$ cells tested, \textbf{1 of 8 rejects}: \cell{red$\cdot$B2} at
$p=0.0004$ against threshold $0.0063$. The next cell, \cell{red$\cdot$B1} at
$p=0.033$ against $0.0071$, does not reject, and the procedure stops.

Two cells (\cell{cls$\cdot$B2}, \cell{cls$\cdot$B1}) have a signal-free pipeline
reproducing at least 80\% of the observed saving --- in \cell{cls$\cdot$B1} the
shuffled median actually exceeds the observed value. One cell
(\cell{red$\cdot$B0}) yields no saving at all.

\begin{figure}[t]
  \centering
  \includegraphics[width=\textwidth]{fig_f12_permutation_control.pdf}
  \caption[Observed saving against the bundle-permutation null]{%
  \textbf{Observed saving against what a signal-free pipeline reproduces.} Per
  cell, the observed SR-lossless cost saving and the null distribution from
  $B=10{,}000$ bundle permutations. The permutation unit is the whole task
  bundle, not the label alone: permuting $y$ by itself disconnects the label
  from the outcomes that define it, and at $B=200$ that error moved two cells in
  opposite directions.}
  \label{fig:f12}
\end{figure}

\section{Why AUROC is not the answer here}
\label{sec:auroc-not-answer}

The one cell that survives, \cell{red$\cdot$B2}, has an AUROC of 0.526 --- below
its own best single covariate at 0.790. That is not a contradiction; the two
quantities measure different things and this data pulls them apart.

AUROC scores the \emph{global} ranking. The saving comes from the \emph{tail}.
The two come apart hardest exactly where the class balance is worst: in a cell
that is 93\% negative, a model can order the sea of true negatives well enough to
score 0.70 without ever confidently identifying a positive, which is the only
thing triage needs. A precision--recall curve would separate these two situations
where AUROC does not, and we report AUROC only because it is the quantity the
routing literature reports; the policy tables in Chapter~\ref{ch:e2e}, not this
column, are what the conclusion rests on.
\cell{red$\cdot$B2} routes 192 of 203 tasks (95\%) to the cheap mode with no
success loss --- in a cell where only 7.4\% of tasks are solvable by anything at
all, so almost nothing in that 95\% was ever going to succeed. It differs from
the free always-cheapest policy on 5\% of the allocation, and those 11 retained
tasks happen to hold 5 of the cell's successes against always-cheapest's 4. The
permutation null is detecting that tail enrichment, not an ordered score.

\begin{scopebox}
\textbf{The honest phrasing.} It is \emph{not} ``the label is predictable, yet
triage fails''. It is: \textbf{at 2--36\% base success, a high AUROC is neither
necessary nor sufficient. What decides whether triage saves anything is whether
a handful of tail tasks land on the right side --- and at $n=203$ that handful
is four successes.}
\end{scopebox}

This has a consequence for how the next chapter must be read. A classifier
metric cannot stand in for a system result in this regime, in either direction:
a cell can discriminate well and buy nothing, and a cell can barely discriminate
and buy something. Chapter~\ref{ch:e2e} therefore evaluates policies, not
scores.

\section{What this chapter establishes}
\label{sec:ch5-summary}

The which-mode half of routing cannot be attempted: 4 of 6 cells have no
trainable classifier, and the constraint is label supply
(Chapter~\ref{ch:disc}).

The triage half --- the easier half by construction --- is weakly predictable,
at AUROC 0.526--0.758, clearing its best single covariate in 5 of 8 cells. But
most of the apparent predictability on VisualWebArena rests on two features no
deployment possesses, and after removing them five of six cells fall. Under a
$B=10{,}000$ bundle-permutation null with Holm correction across the eight
cells, one cell's saving survives, and it is the cell with near-chance global
discrimination.

\chapterhandoff{Chapter~\ref{ch:pred} shows the routing target is at best weakly
predictable under a deployment-faithful information boundary, and that most of
what looked like predictability was a benchmark annotation. Predictive
performance is nonetheless the wrong terminal quantity: a router can rank well
and still lose once false escalations, missed escalations and its own cost are
priced in --- and, as \cell{red$\cdot$B2} shows, the converse also occurs.
Chapter~\ref{ch:e2e} therefore runs the policy end to end and asks whether it
beats the fixed alternative a deployment already has for free.}
